% 2026-11-11 / v33
\section{射影核による低rotation積の相対底変更}\label{gr:sec:basechange}

本節では、最終目標に必要な比較を見直し、一般の非正則基底上のMF核定理を
先に完成させる経路を迂回する。可換側のKuznetsov成分は滑らかな射影族の
admissible部分圏なので、その\emph{真の導来自然変換}はperfectな幾何学的核の
導来Homで計算できる。これと座標切断の持上げを組み合わせると、
次数一・二の全Artin底変更が従う。独立中心変数も有限群対称性も不要である。

To\"enの相対導来Morita定理\cite[Theorems~8.9, 8.15]{GRToen}、
Kuznetsovの射影核によるHochschild理論\cite{GRKuznetsov}、
perfect直接像の底変更\cite[Tag~0B91]{GRStacks}を既知の入力とする。
To\"enの基底$k$は原論文\S2で任意の可換環として固定されており、体に限られない。
最近のLin--Zhang\cite[\S2.1]{GRLZ26}にもprojected kernelを用いる相対的な枠組みがある。
本節はこれらのMorita原理の新規性を主張しない。
体上の低次数同定には\cite[Theorem~5.39]{GRBFK}と
\cite[Theorem~3.5, Remark~3.6, \S5]{GRLZ}だけを用い、
体上の定理を非正則基底へ直接引用することはしない。

\subsection{持ち上がった基底があれば$H^1$の消滅は不要}

\begin{lemma}[閉基底持上げによる$H^0$底変更]\label{gr:lem:lifts}
$(R,\mathfrak m,k)$をArtin局所環、$C\in\operatorname{Perf}(R)$とする。
$H^{-1}(C\otimes_R^{\mathbf L}k)=0$、
$H^0(C\otimes_R^{\mathbf L}k)\simeq k^d$と仮定する。
$\alpha:R^d[0]\to C$が$D(R)$の射で、その剰余の$H^0$が同型なら、
\[
 H^0(\alpha):R^d\xrightarrow{\sim}H^0(C)
 \tag{GR1}\label{gr:eq:lift-isom}
\]
は同型である。任意のArtin商$R\to R'$について自然な写像
\[
 H^0(C)\otimes_R R'\longrightarrow
 H^0(C\otimes_R^{\mathbf L}R')
\]
も同型である。$H^1(C\otimes_R^{\mathbf L}k)$の消滅は仮定しない。
\end{lemma}
\begin{proof}
$C$を有界有限自由複体で表し、各成分の閉代表を選んで$\alpha$を鎖写像とする。
$E=\operatorname{Cone}(\alpha)$も有界有限自由である。
特殊ファイバーの長完全列から$H^{-1}(E_k)=H^0(E_k)=0$。
$E$の有限$\mathfrak m$-進フィルトレーションの各商は
$(\mathfrak m^j/\mathfrak m^{j+1})\otimes_k E_k$なので、
短完全列の長完全列を繰り返すと$H^{-1}(E)=H^0(E)=0$を得る。
$R^d\to C\to E$の長完全列が\eqref{gr:eq:lift-isom}を与える。
同じ鎖写像を$R'$へテンソルし、同じ議論を行うと、両側が
$R'^d$に同定される。これが\emph{自然な}底変更写像と一致する。
\end{proof}

\begin{remark}
$[R\xrightarrow{h}R]$を次数$0,1$に置いた例では、
特殊ファイバーの次数零基底を閉元として持ち上げられない。
上の補題はperfect性や特殊次元だけで底変更を推論しているのではない。
逆に、全基底の閉持上げが既にあれば、次数一の障害を別途計算する必要はない。
\end{remark}

\subsection{相対Kuznetsov成分のprojected kernel}

$R$をArtin局所$\mathbb C$-代数、$N\ge5$、
$X=V(W)\subset\mathbb P_R^{N-1}$を滑らかな相対三次超曲面とする。
$D=\operatorname{Perf}(X)$、$a=N-3$、$E_j=\mathcal O_X(j)$と置き、
\[
 \mathcal A_R=\langle E_0,\ldots,E_{a-1}\rangle^\perp\subset D
\]
をperfect Kuznetsov成分とする。$D^b(\operatorname{coh}X)$全体とは同一視しない。

射影空間の相対線束コホモロジーと
$0\to\mathcal O_{\mathbb P}(j-3)\to\mathcal O_{\mathbb P}(j)\to\mathcal O_X(j)\to0$
により、$E_0,\ldots,E_{a-1}$は$R$-例外列である。
$R\operatorname{Hom}(E_j,E_i)=0$（$j>i$）と
$R\operatorname{Hom}(E_i,E_i)=R$は任意の底変更に両立する。

核の規約を$\Phi_K(M)=Rp_{2*}(p_1^*M\otimes^{\mathbf L}K)$とし、
$K\star L$は$\Phi_K\Phi_L$を表す畳込みとする。
\[
 \mathsf M_j=\operatorname{Cone}(E_j^\vee\boxtimes E_j\longrightarrow\mathcal O_\Delta),
 \qquad P=\mathsf M_0\star\cdots\star\mathsf M_{a-1}
 \tag{GR2}\label{gr:eq:projector}
\]
は順に左mutationと左射影$p=i\ell$を表す。
$i:\mathcal A_R\hookrightarrow D$、$\ell\dashv i$、$\ell i\simeq\mathrm{Id}$である。
右端のmutationから作用することに注意する。

$T=\Delta_*\mathcal O_X(1)$、
\[
 \rho=L_{E_0}\circ(-\otimes\mathcal O_X(1))|_{\mathcal A_R},
 \qquad Q=P\star\mathsf M_0\star T\star P
 \tag{GR3}\label{gr:eq:rotation-kernel}
\]
とする。$M\in\mathcal A_R$なら、$j\ge1$に対し
$R\operatorname{Hom}(E_j,M(1))=R\operatorname{Hom}(E_{j-1},M)=0$であり、
$R\operatorname{Hom}(E_j,E_0)=0$でもある。従って$\rho$は$\mathcal A_R$を保ち、
$Q$は$i\rho\ell$を表す。$Q^{\star r}$は$i\rho^r\ell$を表す。

\begin{theorem}[真の導来Natのperfect模型]\label{gr:thm:perfect-target}
上の状況で、$q:X\times_R X\to\operatorname{Spec}R$として
\[
 \mathcal C_r=Rq_*R\mathcal Hom(P,Q^{\star r})
 \tag{GR4}\label{gr:eq:nat-complex}
\]
は$R$-perfectであり、自然な擬同型
\[
 \mathcal C_r\simeq R\operatorname{Hom}_{\mathcal A_R^e}
       (\Delta_{\mathcal A_R},P_{\rho^r})
\]
を持つ。任意の$R\to R'$に対する導来底変更は、
$X_{R'}$とその同じ射影核から作る複体への自然な擬同型を与える。
比較は自然変換の合成と両立する。
\end{theorem}
\begin{proof}
$X/R$は滑らかなので相対対角は正則埋込みであり、$\mathcal O_\Delta$と$T$はperfect。
従って\eqref{gr:eq:projector}の核はperfectである。
畳込みは滑らかな三重積の射影に沿う導来引戻し・テンソル・proper直接像である。
\cite[Tag~0B91]{GRStacks}により、perfect性と任意の導来底変更との両立が保たれる。
同じ結果を$P^\vee\otimes Q^{\star r}$の直接像へ適用すると、
$\mathcal C_r$のperfect性と底変更が従う。

導来自然変換との同定は体上の次元比較ではない。
\cite[Theorem~8.15]{GRToen}を基底環$R$に適用すると、$D$のenhanced自己関手間の
導来Homは幾何学的核間の導来Homに同定される。
さらに$F\mapsto iF\ell$はenhanced関手圏上で完全忠実である。
実際、$i$上の制限と$\ell i\simeq\mathrm{Id}$から逆向きの写像を得る。
任意の$M\in D$に対する単位$M\to i\ell M$は、
$iF\ell$へ作用させると同型になるため、自然性により
変換は$\mathcal A_R$上の制限から一意に復元される。
同じ単位と三角恒等式はcoherent自然変換複体上の互いに逆な比較を与える。
従って$P$と$Q^{\star r}$の間のHomは、
$\mathrm{Id}_{\mathcal A_R}$と$\rho^r$の間の真の導来自然変換を計算する。

核と変換は全て評価射・mutation・畳込みから作られる。
これらの導来底変更と射の合成は可換であるので、最後の両立性も得る。
無界で項別自由なbar複体をK-flatとみなす段階はない。
\end{proof}

\subsection{座標切断からの全基底持上げ}

任意の$s\in H^0(X,\mathcal O_X(1))$に対し、
\[
 \alpha_s(M):M\xrightarrow{\,s\,}M(1)
       \longrightarrow L_{E_0}(M(1))=\rho(M)
 \tag{GR5}\label{gr:eq:sections}
\]
はenhancedな自然変換である。後半はmutationの単位であり、
核の射としても$R$上で定義される。
$x_1,\ldots,x_N$はこうして$H^0(\mathcal C_1)$の閉持上げを与える。
次数二には$\rho(\alpha_{x_j})\circ\alpha_{x_i}$を用いる。
この段階では順序付き積を使えば十分で、可換性を先取りしない。

\begin{lemma}[特殊ファイバーの低次数入力]\label{gr:lem:field}
$F=W\bmod\mathfrak m$に対し、$r=1,2$では
\[
 H^{-1}(\mathcal C_{r,k})=0,\qquad
 H^0(\mathcal C_{r,k})=\operatorname{Jac}(F)_r.
 \tag{GR6}\label{gr:eq:field}
\]
\eqref{gr:eq:sections}の特殊化はJacobian環の$x_i$に対応し、
その合成は多項式積に対応する。従って次数一の$N$個の座標と、
次数二で適当に選んだ$\binom N2$個の二次単項式が基底を与える。
\end{lemma}
\begin{proof}
体上のBFK公式とLin--Zhangの乗法同定を使う
\cite[Theorem~5.39]{GRBFK}\cite[Theorem~3.5, Remark~3.6, \S5]{GRLZ}。
恒等sectorは偶数コホモロジー次数$j=2b$に$\operatorname{Jac}(F)_{r+3b}$を与える。
非恒等sectorは、$j\equiv N\pmod2$かつ
\[
 2r+3j-N=0
\]
のときだけ寄与する。$j=-1$なら$N=2r-3\le1$、
$j=0$なら$N=2r\le4$なので、$N\ge5$ではいずれも起こらない。
特殊ファイバーの勾配は$N$個の二次正則列なので、
次数一・二の次元は$N,\binom N2$である。
座標による変換とその合成は同公式の恒等sectorの多項式作用である。
ここではrotation自体を入力するため$\gcd(N,3)=1$を要求しない。
\end{proof}

\begin{theorem}[全Artin低rotation底変更]\label{gr:thm:bc}
任意のArtin局所$\mathbb C$-代数$R$と、滑らかな
$X=V(W)\subset\mathbb P_R^{N-1}$（$N\ge5$）について
\[
 \operatorname{Nat}_{dg}^0(\mathrm{Id},\rho)\simeq R^N,
 \qquad
 \operatorname{Nat}_{dg}^0(\mathrm{Id},\rho^2)\simeq R^{\binom N2}.
 \tag{GR7}\label{gr:eq:ranks}
\]
自然な$H^0$の底変更は全Artin商で同型で、合成と両立する。
独立な中心変数・有限群対称性・$H^1$の消滅は仮定しない。
\end{theorem}
\begin{proof}
定理\ref{gr:thm:perfect-target}により真の標的をperfect複体$\mathcal C_r$で計算できる。
$r=1$では$\alpha_{x_i}$を用い、$r=2$では特殊ファイバーの基底となる
$\binom N2$個の順序付き合成を選ぶ。
これらはそれぞれ$R^N[0]\to\mathcal C_1$、
$R^{\binom N2}[0]\to\mathcal C_2$という射を与える。
補題\ref{gr:lem:field}と補題\ref{gr:lem:lifts}を適用する。
全ての射が核上で底変更と可換なので、得られる$H^0$の同型も自然なものであり、
積と底変更の両立が従う。
\end{proof}

\begin{remark}[障害群を消すより短い理由]
五変数の$r=1$では特殊ファイバーの$H^1$は二次元であり、
七変数の$r=2$でも二つの非恒等sectorが$H^1$に寄与する。
それでも上の定理は成立する。幾何学的座標によって必要な全基底が
既に閉元として持ち上がっているからであり、障害群そのものを零にする必要はない。
これは以前の中心対称性による消滅結果を否定するものではなく、
本問題の可換標的には不要な追加条件を取り除くものである。
\end{remark}

\subsection{既存Artin比較に使う生成性の仮定の点検}

\begin{lemma}[導来特殊化とcompact生成子]\label{gr:lem:derived-generation}
$R$をArtin局所$\mathbb C$-代数、$\mathcal B$を小さな冪等完備pretriangulated
$R$-線形dg圏とする。全Hom複体が$R$-K-flatであり、通常の係数特殊化が
導来Morita特殊化$\mathcal B\otimes_R^{\mathbf L}k$を表すとする。
$G\in\mathcal B$の剰余がその導来特殊ファイバーを生成すれば、$G$は$\mathcal B$を生成する。
とくに、\emph{項別平坦性だけ}をこの仮定の代わりに使うことはできない。
\end{lemma}
\begin{proof}
$\mathcal D=\operatorname{Ind}(\mathcal B)$を取る。$G$はcompactであり、
$R\operatorname{Hom}(G,-)$は全colimitと可換する。
compact対象のHomの導来底変更は表現可能対象で成立し、
両辺がcolimitと可換することから任意の$M\in\mathcal D$へ延長される。
$R\operatorname{Hom}(G,M)=0$なら、特殊ファイバーの生成性から
$M\otimes_R^{\mathbf L}k=0$。
$R$は$D(R)$内で$k$の有限extensionから作られるので、
$M=M\otimes_R^{\mathbf L}R=0$である。
従って$G$は$\mathcal D$のcompact generatorであり、compact部分を取ると結論を得る。
\end{proof}

本研究の正の内部次数付き有限rank MFでは、固定cohomological次数のHomは
有限自由$R_m$-加群であり、特殊ファイバーのHomの全cohomologyは有限次元かつ有界である。
特殊Hom複体の体上の分裂収縮を選ぶと、$h$の冪零性により
homological perturbationの逆は有限和となる。
従って実際のHom複体は有界有限自由$R_m$-複体へホモトピー同値であり、K-flatである。
これは既存Artin残余比較の\texttt{am:lem:mf-hom}の入力である。
特殊化関手はHomで導来底変更を表し、その像が古典生成子を含むことを
windowの三つのstabilizationで確認してから、上の補題を適用できる。
零曲率の任意の二周期圏で項別自由性だけを使う旧一般補題の反例とは矛盾しない。
本節ではこの必要なK-flat性・導来特殊化・生成子の像を分離しており、
形式MF復元やstabilizationの全証明を再証明したとはしない。

\subsection{非可換入力への適用と依存する既存結果}

以下で用いる既存入力を明記する。
\begin{description}
\item[(F)] 全有限Artin切断で、指定rotationと底変更を保つ
$\mathcal K^{\rm lin}_{B_m}\simeq\operatorname{Ku}_{\rm perf}(Y_m/R_m)$がある。
これは現行論文Bの形式MF復元とArtin残余比較を合成した結果である。
本節の幾何学的定理は(F)に依存しないが、非可換側への適用は依存する。
\item[(D)] 計算した有限対角の低次数$H^0$は自由で、全Artin底変更に両立し、
実際の残余pairへの乗法的characteristic写像$\Theta$がある。
これは局所正則対角、乗法的bar比較、および(F)のwindowとの両立を入力とする。
\end{description}
ここに一般のMF kernel--functor比較(KF)を追加してはならない。
(F)と(D)は既存証明への依存であり、本節だけで全既存証明を独立再監査したとの意味ではない。

\begin{corollary}[既存形式比較の下での(BC)の充足]\label{gr:cor:nc}
(F)が成立する中心三次対について、標的側の条件(BC)は
定理\ref{gr:thm:bc}から従う。(D)も成立し、特殊ファイバーの$\Theta_1,\Theta_2$が
Jacobianの標準同定なら、両写像は全有限Artin切断で同型である。
\end{corollary}
\begin{proof}
(F)をdg双加群の導来自然変換へ適用し、標的を\eqref{gr:eq:nat-complex}で計算する。
Morita同値はrotation双加群とその合成を保つ。
標的は定理\ref{gr:thm:bc}により所定の階数の自由加群で底変更に両立する。
(D)のsourceも同じ階数の自由加群である。
$\Theta_r$の行列式は剰余で非零なので単元であり、同型である。
乗法性と全切断の整合性は入力した$\Theta$のものを保つ。
\end{proof}

可換側で低積を厳密にJacobian積と同定するには、既存の相対window比較を用いる。
可換curved Hochschild複体では導分$\partial_i$の微分が、符号を除いて$t\partial_iW$である。
従って$t^{-1}\partial_i$は$\partial_iW$の有限arityの原像を与える。
characteristic写像により、偏微分による自然変換は境界になる。
windowは多項式倍と内部shiftを保ち、mutationの単位は商で同型になるので、
これは\eqref{gr:eq:sections}の多項式作用と一致する。
相対勾配Koszulの完全性から$\operatorname{Jac}(W)_1,\operatorname{Jac}(W)_2$は所定の階数で自由。
これらから$H^0(\mathcal C_r)$への写像は特殊化で同型なので、全Artinで同型である。
ここでもMF核の全忠実性は用いない。

\begin{corollary}[計算済み三次式の形式的同定]\label{gr:cor:formal-polynomial}
既存入力(F),(D)を満たす$\lambda=1$の混合族について、
有限対角から得た多項式$G_h+v^3$と形式復元先$W_h$は、
$h=0$で恒等な$\operatorname{GL}_5(\mathbb C[[h]])$変換で同型である。
従って全ての$m$で、非可換線束残余pairは
$V(G_h+v^3\bmod h^{m+1})$のperfect Kuznetsov pairと同値である。
これは非零複素パラメータのgeneric圏同値を主張しない。
\end{corollary}
\begin{proof}
系\ref{gr:cor:nc}により低次数の積が同定され、次数二の核として勾配空間が一致する。
次数一の基底の対応は特殊ファイバーで恒等に規格化でき、
全切断の互換性により逆極限の線形対応を得る。
この座標変更後、$\nabla W_h=A_h\nabla(G_h+v^3)$、$A_h\equiv1\pmod h$と書ける。
三次式の対称三重形式を$T$とすると、$A_h^{\mathsf t}$は
$T(A_h^{\mathsf t}u,v,w)=T(u,A_h^{\mathsf t}v,w)=T(u,v,A_h^{\mathsf t}w)$を満たす。
$A_h^{\mathsf t}$の形式三乗根も同じ性質を持つため、
$C_h=(A_h^{\mathsf t})^{1/3}$によって
$(G_h+v^3)(C_hx)=W_h(x)$となる。
この形式座標同型と(F)を各切断で合成する。
\end{proof}

\subsection{generic復元までの最短経路と残る証明}

本節により、既存の形式残余比較を利用できる範囲では、
一般の相対MF核比較(KF)、無界barの全面的な特殊化理論、
中心対称性による次数一障害の逐次処理を、最終目標の前提にする必要がなくなった。
有限代数的抽出という目標は維持し、標的の処理だけを幾何学的射影核へ置き換える。

しかし\emph{式の代数性}と\emph{圏同値の代数性}は別である。
(F)から得た全Artin同値や上の形式座標変換から、
非可換側の指定生成子がgenericでも生成するとは推論できない。
次に必要なのは、代数的な残余模型と$\operatorname{Ku}(V(G_h+v^3))$を結ぶ
双加群、逆双加群、両合成の単位・余単位、rotation intertwinerを
\emph{共通の有限cell/retractデータ}で構成することである。
その有限データがあれば、構造方程式・逆写像・ホモトピーの恒等式を
有限の代数方程式として検証し、存在する開集合へ広げられる。
現時点でその有限証明書、一般混合型の相対対角perfect性、generic生成性は未構成である。
滑らかなproper dg代数の有限$A_n$決定性\cite{GRToenFinite}を利用する場合も、
非可換側のsmoothnessを先に確認する必要があり、ここを循環させない。

本節の計算コードは、体上sector重みと閉基底持上げ補題の有限自由複体例を検査する。
幾何学的核定理、既存の形式MF同値、または非可換対角の全行列を計算機で検証したものではない。
